% GENERATED FILE. Edit canonical lessons, then run npm run generate:books.
% canonical-source-hash: fnv1a64:efa79bfc17b04261
% canonical-lessons: UR-C447-itminan, UR-C447-asuda, UR-C447-qarar, UR-C447-dilasa, UR-C447-khushgavar, UR-R447-review-440-447

\chapter{Calm Again}
\label{ch:ur-calm-again}
\hlchaptermodality{447}

\begin{chapteropening}
\textbf{By the end of this chapter:} \emph{I can say that someone felt calm and content.}

Complete the last lesson of chapter 447: I can say that someone felt calm and content.
\end{chapteropening}

\section[\texorpdfstring{\ur{اطمینان} (itmīnān)}{اطمینان (itmīnān)}]{\ur{اطمینان} (itmīnān) — satisfaction; peace of mind}

\label{lesson:UR-C447-itminan}



\begin{quote}
Before the new one: say the Urdu for to writhe, then the Urdu for to throb.
\end{quote}

\subsection*{You'll want to know: \ur{اطمینان}}
\textbf{\ur{اطمینان}} — \emph{itmīnān} — \textquotedblleft{}satisfaction; peace of mind\textquotedblright{}.

\begin{quote}
\ur{کام} \ur{ختم} \ur{کر} \ur{کے} \ur{مجھے} \ur{اطمینان} \ur{ہوا۔}
\end{quote}

\begin{quote}
\emph{kām khatm kar ke mujhe itmīnān huā.}
\end{quote}

\begin{quote}
I felt satisfied once the work was finished.
\end{quote}

Masculine.

\begin{cousinweb}
\textbf{\ur{اطمینان}} (\emph{itmīnān}) is usually traced, through Persian, to Arabic \emph{iṭmiʾnān}, \textquotedblleft{}being at rest\textquotedblright{}.
\end{cousinweb}

\begin{grammarlens}[title={What you've built}]
One more word for describing what you lived and felt.
\end{grammarlens}

\subsection*{Your turn}
\begin{itemize}
\raggedright
  \item \emph{Say it:} \emph{itmīnān}
  \item \emph{Say it:} \emph{itmīnān}, once more
  \item \emph{Say it:} \emph{kām khatm kar ke mujhe itmīnān huā.} — I felt satisfied once the work was finished.
  \item \emph{Recall:} say the Urdu for to writhe, then the Urdu for to throb, then say \emph{itmīnān} again
\end{itemize}

\begin{tcolorbox}[breakable,colback=teal!4,colframe=teal!35!black,title={Before you move on}]
What does \ur{اطمینان} mean? (\textquotedblleft{}Satisfaction; peace of mind\textquotedblright{}.) Say it once more.
\end{tcolorbox}
\section[\texorpdfstring{\ur{آسودہ} (āsūdā)}{آسودہ (āsūdā)}]{\ur{آسودہ} (āsūdā) — at ease; content}

\label{lesson:UR-C447-asuda}



\begin{quote}
Before the new one: say the Urdu for to throb, then the Urdu for satisfaction.
\end{quote}

\subsection*{You'll want to know: \ur{آسودہ}}
\textbf{\ur{آسودہ}} — \emph{āsūdā} — \textquotedblleft{}at ease; content\textquotedblright{}.

\begin{quote}
\ur{اچھا} \ur{کھانا} \ur{کھا} \ur{کر} \ur{سب} \ur{آسودہ} \ur{تھے۔}
\end{quote}

\begin{quote}
\emph{achchhā khānā khā kar sab āsūdā the.}
\end{quote}

\begin{quote}
After a good meal everyone was content.
\end{quote}

\begin{cousinweb}
\textbf{\ur{آسودہ}} (\emph{āsūdā}) is usually traced to Persian \emph{āsūda}, \textquotedblleft{}rested, at ease\textquotedblright{}.
\end{cousinweb}

\begin{grammarlens}[title={What you've built}]
One more word for describing what you lived and felt.
\end{grammarlens}

\subsection*{Your turn}
\begin{itemize}
\raggedright
  \item \emph{Say it:} \emph{āsūdā}
  \item \emph{Say it:} \emph{āsūdā}, once more
  \item \emph{Say it:} \emph{achchhā khānā khā kar sab āsūdā the.} — After a good meal everyone was content.
  \item \emph{Recall:} say the Urdu for to throb, then the Urdu for satisfaction, then say \emph{āsūdā} again
\end{itemize}

\begin{tcolorbox}[breakable,colback=teal!4,colframe=teal!35!black,title={Before you move on}]
What does \ur{آسودہ} mean? (\textquotedblleft{}At ease; content\textquotedblright{}.) Say it once more.
\end{tcolorbox}
\section[\texorpdfstring{\ur{قرار} (qarār)}{قرار (qarār)}]{\ur{قرار} (qarār) — composure; rest (of the mind)}

\label{lesson:UR-C447-qarar}



\begin{quote}
Before the new one: say the Urdu for satisfaction, then the Urdu for at ease.
\end{quote}

\subsection*{You'll want to know: \ur{قرار}}
\textbf{\ur{قرار}} — \emph{qarār} — \textquotedblleft{}composure; rest (of the mind)\textquotedblright{}.

\begin{quote}
\ur{بیٹے} \ur{کی} \ur{خبر} \ur{سن} \ur{کر} \ur{ماں} \ur{کو} \ur{قرار} \ur{آیا۔}
\end{quote}

\begin{quote}
\emph{beṭe kī khabar sun kar māṅ ko qarār āyā.}
\end{quote}

\begin{quote}
Hearing news of her son, the mother was at peace.
\end{quote}

Masculine.

\begin{cousinweb}
\textbf{\ur{قرار}} (\emph{qarār}) is usually traced, through Persian, to Arabic \emph{qarār}, \textquotedblleft{}settling, rest\textquotedblright{}.
\end{cousinweb}

\begin{grammarlens}[title={What you've built}]
One more word for describing what you lived and felt.
\end{grammarlens}

\subsection*{Your turn}
\begin{itemize}
\raggedright
  \item \emph{Say it:} \emph{qarār}
  \item \emph{Say it:} \emph{qarār}, once more
  \item \emph{Say it:} \emph{beṭe kī khabar sun kar māṅ ko qarār āyā.} — Hearing news of her son, the mother was at peace.
  \item \emph{Recall:} say the Urdu for satisfaction, then the Urdu for at ease, then say \emph{qarār} again
\end{itemize}

\begin{tcolorbox}[breakable,colback=teal!4,colframe=teal!35!black,title={Before you move on}]
What does \ur{قرار} mean? (\textquotedblleft{}Composure; rest (of the mind)\textquotedblright{}.) Say it once more.
\end{tcolorbox}
\section[\texorpdfstring{\ur{دلاسا} (dilāsā)}{دلاسا (dilāsā)}]{\ur{دلاسا} (dilāsā) — consolation; comforting words}

\label{lesson:UR-C447-dilasa}



\begin{quote}
Before the new one: say the Urdu for at ease, then the Urdu for composure.
\end{quote}

\subsection*{You'll want to know: \ur{دلاسا}}
\textbf{\ur{دلاسا}} — \emph{dilāsā} — \textquotedblleft{}consolation; comforting words\textquotedblright{}.

\begin{quote}
\ur{دوست} \ur{نے} \ur{مجھے} \ur{دلاسا} \ur{دیا۔}
\end{quote}

\begin{quote}
\emph{dost ne mujhe dilāsā diyā.}
\end{quote}

\begin{quote}
My friend comforted me.
\end{quote}

Masculine.

\begin{cousinweb}
\textbf{\ur{دلاسا}} (\emph{dilāsā}) is usually traced to Persian \emph{dilāsā}, \textquotedblleft{}comfort\textquotedblright{}, a word built on \emph{dil}, \textquotedblleft{}heart\textquotedblright{} (\textbf{\ur{دل}}).
\end{cousinweb}

\begin{grammarlens}[title={What you've built}]
One more word for describing what you lived and felt.
\end{grammarlens}

\subsection*{Your turn}
\begin{itemize}
\raggedright
  \item \emph{Say it:} \emph{dilāsā}
  \item \emph{Say it:} \emph{dilāsā}, once more
  \item \emph{Say it:} \emph{dost ne mujhe dilāsā diyā.} — My friend comforted me.
  \item \emph{Recall:} say the Urdu for at ease, then the Urdu for composure, then say \emph{dilāsā} again
\end{itemize}

\begin{tcolorbox}[breakable,colback=teal!4,colframe=teal!35!black,title={Before you move on}]
What does \ur{دلاسا} mean? (\textquotedblleft{}Consolation; comforting words\textquotedblright{}.) Say it once more.
\end{tcolorbox}
\section[\texorpdfstring{\ur{خوشگوار} (khushgavār)}{خوشگوار (khushgavār)}]{\ur{خوشگوار} (khushgavār) — pleasant}

\label{lesson:UR-C447-khushgavar}



\begin{quote}
Before the new one: say the Urdu for composure, then the Urdu for consolation.
\end{quote}

\subsection*{You'll want to know: \ur{خوشگوار}}
\textbf{\ur{خوشگوار}} — \emph{khushgavār} — \textquotedblleft{}pleasant\textquotedblright{}.

\begin{quote}
\ur{آج} \ur{کا} \ur{موسم} \ur{بڑا} \ur{خوشگوار} \ur{ہے۔}
\end{quote}

\begin{quote}
\emph{āj kā mausam baṛā khushgavār hai.}
\end{quote}

\begin{quote}
The weather is very pleasant today.
\end{quote}

\begin{cousinweb}
\textbf{\ur{خوشگوار}} (\emph{khushgavār}) is usually traced to Persian \emph{khushgawār}, \textquotedblleft{}agreeable\textquotedblright{}; its first part is \textbf{\ur{خوش}}, \textquotedblleft{}happy\textquotedblright{}.
\end{cousinweb}

\begin{grammarlens}[title={What you've built}]
One more word for describing what you lived and felt.
\end{grammarlens}

\subsection*{Your turn}
\begin{itemize}
\raggedright
  \item \emph{Say it:} \emph{khushgavār}
  \item \emph{Say it:} \emph{khushgavār}, once more
  \item \emph{Say it:} \emph{āj kā mausam baṛā khushgavār hai.} — The weather is very pleasant today.
  \item \emph{Recall:} say the Urdu for composure, then the Urdu for consolation, then say \emph{khushgavār} again
\end{itemize}

\begin{tcolorbox}[breakable,colback=teal!4,colframe=teal!35!black,title={Before you move on}]
What does \ur{خوشگوار} mean? (\textquotedblleft{}Pleasant\textquotedblright{}.) Say it once more.
\end{tcolorbox}
\section[(dialogue)]{Looking back: from Naming a Feeling to Calm Again}

\label{lesson:UR-R447-review-440-447}



\begin{quote}
Say the Urdu for an emotion and a feeling.
\end{quote}

\subsection*{Your turn}
Say these aloud:

\begin{itemize}
\raggedright
  \item an emotion, a feeling, to sense, a state of mind, a mood
  \item delight, enjoyment, joyful, brimming over, a surge of eagerness
  \item astonishment, strange, dumbfounded, a wonder, an amazing thing
  \item grief, sorrow, disappointed, to fret inwardly, dejected
  \item anger, to get irritated, incensed, displeased, a fit of rage
  \item dread, to frighten, to cower, to crouch out of sight, to huddle
  \item to worry, anxiety, a burden, to writhe, to throb
  \item satisfaction, at ease, composure, consolation, pleasant
\end{itemize}

\begin{tcolorbox}[breakable,colback=teal!4,colframe=teal!35!black,title={Before you move on}]
An emotion? (\textbf{\ur{جذبہ}}, \emph{jazbā}.) A feeling? (\textbf{\ur{احساس}}, \emph{ehsās}.) To sense? (\textbf{\ur{محسوس} \ur{کرنا}}, \emph{mahsūs karnā}.) A state of mind? (\textbf{\ur{کیفیت}}, \emph{kaifiyat}.) A mood? (\textbf{\ur{مزاج}}, \emph{mizāj}.) Delight? (\textbf{\ur{مسرت}}, \emph{masarrat}.) Enjoyment? (\textbf{\ur{لطف}}, \emph{lutf}.) Joyful? (\textbf{\ur{شادمان}}, \emph{shādmān}.) Brimming over? (\textbf{\ur{سرشار}}, \emph{sarshār}.) A surge of eagerness? (\textbf{\ur{امنگ}}, \emph{umang}.) Astonishment? (\textbf{\ur{تعجب}}, \emph{tājjub}.) Strange? (\textbf{\ur{عجیب}}, \emph{ajīb}.) Dumbfounded? (\textbf{\ur{دنگ}}, \emph{dang}.) A wonder? (\textbf{\ur{عجوبہ}}, \emph{ajūbā}.) An amazing thing? (\textbf{\ur{اچنبھا}}, \emph{achambhā}.) Grief? (\textbf{\ur{غم}}, \emph{gham}.) Sorrow? (\textbf{\ur{دکھ}}, \emph{dukh}.) Disappointed? (\textbf{\ur{مایوس}}, \emph{māyūs}.) To fret inwardly? (\textbf{\ur{کڑھنا}}, \emph{kuṛhnā}.) Dejected? (\textbf{\ur{افسردہ}}, \emph{afsurdā}.) Anger? (\textbf{\ur{غصہ}}, \emph{ghussā}.) To get irritated? (\textbf{\ur{چڑنا}}, \emph{chiṛnā}.) Incensed? (\textbf{\ur{برہم}}, \emph{barham}.) Displeased? (\textbf{\ur{خفا}}, \emph{khafā}.) A fit of rage? (\textbf{\ur{طیش}}, \emph{taish}.) Dread? (\textbf{\ur{خوف}}, \emph{khauf}.) To frighten? (\textbf{\ur{ڈرانا}}, \emph{ḍarānā}.) To cower? (\textbf{\ur{سہمنا}}, \emph{sahamnā}.) To crouch out of sight? (\textbf{\ur{دبکنا}}, \emph{dubaknā}.) To huddle? (\textbf{\ur{سمٹنا}}, \emph{simaṭnā}.) To worry? (\textbf{\ur{فکر} \ur{کرنا}}, \emph{fikr karnā}.) Anxiety? (\textbf{\ur{تشویش}}, \emph{tashvīsh}.) A burden? (\textbf{\ur{بوجھ}}, \emph{bojh}.) To writhe? (\textbf{\ur{تڑپنا}}, \emph{taṛapnā}.) To throb? (\textbf{\ur{دھڑکنا}}, \emph{dhaṛaknā}.) Satisfaction? (\textbf{\ur{اطمینان}}, \emph{itmīnān}.) At ease? (\textbf{\ur{آسودہ}}, \emph{āsūdā}.) Composure? (\textbf{\ur{قرار}}, \emph{qarār}.) Consolation? (\textbf{\ur{دلاسا}}, \emph{dilāsā}.) Pleasant? (\textbf{\ur{خوشگوار}}, \emph{khushgavār}.)
\end{tcolorbox}
